%%%%%%%%%%  scan idx 24  |  folio 11  %%%%%%%%%%
\[
T_f\in K^\bullet(X).
\]
Lorsque $f$ admet une factorisation (2.1) par une immersion dans un sch\'ema relatif lisse, la d\'efinition de $T_f$ n'offre pas de nouvelle difficult\'e. On peut d'ailleurs la reformuler dans ce cas, d'une fa\c{c}on plus conforme \`a l'esprit «~cat\'egories d\'eriv\'ees~», en consid\'erant l'homomorphisme canonique bien connu
\begin{equation}\tag{4.1}
\underline N_{X/X'}\xrightarrow{d}\Omega^1_{X'/Y}\otimes_{\underline O_{X'}}\underline O_X
\end{equation}
induit par la diff\'erentielle ext\'erieure
\[
d_{X'/X}:\underline J\subset\underline O_{X'}\longrightarrow\Omega^1_{X'/Y},
\]
et le complexe de chaines
\begin{equation}\tag{4.2}
L^{X/Y}_\bullet\in\operatorname{Ob}D(X)
\end{equation}
d\'efini gr\^ace \`a (4.1) «~en prolongeant par des z\'eros~», complexe \'egal en degr\'e $1$ resp.\ $0$ \`a la source resp.\ au but de l'homomorphisme (4.1). La formule (2.2) prend alors la forme plus sympathique
\begin{equation}\tag{4.3}
\check T_f=\cl^\bullet\bigl(L^{X/Y}_\bullet\bigr),
\end{equation}
formule qui a un sens dans $K^\bullet(X)$ lorsque $f$ est un morphisme localement d'intersection compl\`ete, puisque dans ce cas le complexe $L^{X/Y}_\bullet$ est manifestement parfait. L'assertion d'invariance de $T_f$ relativement \`a la factorisation choisie de $f$ peut alors se pr\'eciser facilement en une assertion d'invariance pour le complexe $L^{X/Y}_\bullet$~: modulo isomorphisme canonique dans la cat\'egorie d\'eriv\'ee $D(X)$, ce complexe est \'egalement ind\'ependant de la factorisation choisie (Exp.\ VIII).

\subsubsection*{4.4.}
Lorsque $f:X\to Y$ est un morphisme quelconque de sch\'emas, on peut encore construire un objet canonique $L^{X/Y}_\bullet$ de la cat\'egorie d\'eriv\'ee $D(X)$, d\'efini \`a isomorphisme unique pr\`es, qui sur chaque ouvert affine de $X$ au dessus d'un ouvert affine de $Y$ soit isomorphe au complexe construit suivant le principe expos\'e \`a l'alin\'ea pr\'ec\'edent (o\`u il convient cependant dans le cas g\'en\'eral de remplacer «~lisse~» par «~formellement lisse~», de sorte que toute alg\`ebre $B$ sur un anneau $A$ apparait comme quotient d'une alg\`ebre formellement lisse, par exemple une alg\`ebre de polyn\^omes). Pour une telle d\'efinition il n'est pas possible de proc\'eder par simple «~recollement~» \`a partir des constructions affines, car les objets de la cat\'egorie d\'eriv\'ee $D(X)$ sont de nature
